% =============================================================================
% KAPITEL 4: METHODIK
% =============================================================================
% Dieses Kapitel beschreibt das experimentelle Design der Arbeit: wie das
% Experiment durchgeführt wird und anhand welcher Kriterien Erfolg bzw.
% Misserfolg des Ansatzes bewertet wird. Die konkrete Implementierung --
% Datengrundlage, Feature Engineering und Umsetzung -- folgt in Kapitel 5.
% =============================================================================

\chapter{Methodik}\label{chap:methodik}

Dieses Kapitel legt das experimentelle Design der vorliegenden Arbeit dar.
Abschnitt~\ref{sec:forschungsdesign} beschreibt den übergeordneten Ansatz
und die Experimentstruktur. Abschnitt~\ref{sec:split} erläutert die
Datenaufteilung und das Vorgehen zur Vermeidung von Data Leakage.
Abschnitt~\ref{sec:metriken} definiert die Evaluationsmetriken, anhand
derer Erfolg und Misserfolg des Ansatzes beurteilt werden.
Abschnitt~\ref{sec:methodik-zusammenfassung} fasst die methodischen
Entscheidungen zusammen.

% -----------------------------------------------------------------------------
\section{Forschungsdesign und Experimentstruktur}
\label{sec:forschungsdesign}
% -----------------------------------------------------------------------------

Das zentrale Forschungsziel ist ein \textbf{kontrollierter Vergleich} zwischen
zwei Feature-Sets: einem Basis-Set ohne Expected-Goals-Daten und einem
erweiterten Set mit xG-Metriken. Alle anderen Bedingungen -- Algorithmen,
Hyperparameter-Suchraum, Datensplit und Evaluationsmetriken -- bleiben
identisch, sodass beobachtete Leistungsunterschiede kausal auf die
Integration der xG-Features zurückgeführt werden können.

\subsection{Zwei-Gruppen-Design}

Jeder Modelllauf wird zweimal durchgeführt: einmal mit dem Basis-Feature-Set
(25 Features, ohne xG) und einmal mit dem erweiterten Feature-Set (34
Features, inklusive xG). Der direkte Paarvergleich dieser Gruppen bildet
die empirische Grundlage für die Beantwortung der Forschungsfrage
\cite{JamesEtAl2021}.

\subsection{Datensatzvarianten}

Um die Robustheit der Ergebnisse zu gewährleisten, werden für jede der
beiden Feature-Gruppen weitere Vorverarbeitungsvarianten gebildet, die
in Kapitel~\ref{chap:modellierung} im Detail beschrieben werden.
Dadurch entstehen insgesamt \textbf{zwölf Datensatzvarianten}, die sich
entlang dreier Achsen unterscheiden: xG-Integration (mit/ohne),
Ausreißerbehandlung (unbereinigt/IQR-Bereinigung) und Dimensionsreduktion
(alle Features/Feature-Selektion/PCA). Diese Struktur erlaubt es, den
Effekt jeder Vorverarbeitungsentscheidung isoliert zu untersuchen
\cite{BunkerThabtah2019}.

\subsection{Algorithmenauswahl}

Auf jeder der zwölf Varianten werden fünf Klassifikationsalgorithmen
trainiert: Logistische Regression (Baseline), Random Forest, LightGBM,
XGBoost und ein Multilayer Perceptron. Die Auswahl deckt ein breites
Spektrum von linearen Modellen über Ensemble-Methoden bis zu neuronalen
Netzen ab und ermöglicht eine algorithmisch robuste Einschätzung des
xG-Effekts \cite{BunkerThabtah2019,GrollLeyEffenberger2019}. Für jeden
Algorithmus wird ein Hyperparameter-Tuning via \texttt{RandomizedSearchCV}
mit \texttt{StratifiedKFold} (5~Folds) durchgeführt, sodass alle Modelle
unter vergleichbaren, optimierten Bedingungen evaluiert werden
\cite{JamesEtAl2021}.

% -----------------------------------------------------------------------------
\section{Datenaufteilung und Leakage-Prävention}
\label{sec:split}
% -----------------------------------------------------------------------------

Der Datensatz wird im Verhältnis \textbf{80:20} in einen Trainings- und
einen Testdatensatz aufgeteilt. Durch Stratifizierung bleibt die
Klassenverteilung (Heimsieg, Unentschieden, Auswärtssieg) in beiden
Teilmengen repräsentativ erhalten \cite{JamesEtAl2021}. Das Testset wird
ausschließlich für die abschließende Evaluation verwendet und fließt zu
keinem Zeitpunkt in das Training oder das Hyperparameter-Tuning ein.

Data Leakage -- das unbeabsichtigte Einfließen von Zukunftsinformationen
in das Modell -- ist bei Sportvorhersagen ein bekanntes methodisches
Problem \cite{KaufmanRossetPerlich2012}. Diese Arbeit adressiert es auf
zwei Ebenen: Erstens werden bei der Berechnung der Form-Features eines
Spiels ausschließlich \emph{vorangegangene} Spiele berücksichtigt.
Zweitens werden xG-Saisondaten nach einer Lag-1-Strategie integriert:
Für abgeschlossene Saisons dienen die xG-Werte der Vorsaison als Prior,
da Saisonendstatistiken erst nach dem letzten Spieltag feststehen.
Beide Maßnahmen stellen sicher, dass das Modell bei jeder Vorhersage
ausschließlich auf Information zurückgreift, die zum Zeitpunkt des Spiels
tatsächlich verfügbar wäre \cite{KaufmanRossetPerlich2012}.

% -----------------------------------------------------------------------------
\section{Evaluationsmetriken}
\label{sec:metriken}
% -----------------------------------------------------------------------------

Die Wahl der Evaluationsmetriken orientiert sich an der spezifischen
Herausforderung der Fußballvorhersage: drei unausgewogene Klassen, bei
denen Unentschieden systematisch schwerer vorherzusagen sind
\cite{HeGarcia2009}.

\subsection{Primärmetrik: F1 macro}

Die \textbf{Hauptmetrik} ist der \textbf{F1 macro}, das arithmetische
Mittel der klassenspezifischen F1-Scores \cite{Powers2020}:
\begin{equation}
    F1_{\text{macro}} = \frac{1}{K} \sum_{k=1}^{K}
    \frac{2 \cdot \text{Precision}_k \cdot \text{Recall}_k}
         {\text{Precision}_k + \text{Recall}_k}
\end{equation}
mit $K = 3$ Klassen. Anders als die Accuracy, die durch die häufige
Klasse (Heimsieg, \textasciitilde 45\,\%) dominiert wird, behandelt
F1 macro alle Klassen gleich und ist damit robust gegenüber der
Klassen-Imbalance \cite{Powers2020,HeGarcia2009}. Ein Modell, das
Unentschieden vollständig ignoriert, wird durch diese Metrik deutlich
schlechter bewertet als durch die reine Accuracy.

\subsection{Sekundärmetriken}

Ergänzend werden folgende Metriken erhoben:

\begin{itemize}
    \item \textbf{F1 Draw:} Der klassenspezifische F1-Score der
          Klasse \texttt{D}, separat ausgewiesen als zentraler
          Indikator für die Qualität der Unentschieden-Vorhersage.
    \item \textbf{Accuracy:} Anteil korrekt klassifizierter Spiele,
          als verbreitete Baseline-Metrik für den Literaturvergleich
          \cite{BunkerThabtah2019,HubacekSourekZelezny2019}.
    \item \textbf{Precision macro} und \textbf{Recall macro:}
          Differenzieren, ob Fehler eher in falsch positiven oder
          falsch negativen Vorhersagen liegen.
    \item \textbf{Konfusionsmatrix:} Visualisiert die Verteilung
          richtiger und falscher Vorhersagen je Klasse und macht
          systematische Schwächen sichtbar.
\end{itemize}

\subsection{Erfolgskriterium}

Ein Ansatz gilt als \textbf{erfolgreich} im Sinne der Forschungsfrage,
wenn das xG-erweiterte Feature-Set gegenüber dem Basis-Set einen
konsistenten Vorteil im F1 macro über mehrere Algorithmen und Varianten
hinweg zeigt. Ein isolierter Ausreißer in einer einzigen Konfiguration
wird nicht als belastbarer Beleg gewertet. Als \textbf{Misserfolg} gilt,
wenn xG-Features keinen oder einen negativen Effekt haben -- was
ebenfalls ein wissenschaftlich wertvolles Ergebnis darstellt, da es
die Grenzen der Metrik im Vorhersagekontext aufzeigen würde
\cite{AnzerBauer2021,VanEetveldeEtAl2024}.

% -----------------------------------------------------------------------------
\section{Zusammenfassung}
\label{sec:methodik-zusammenfassung}
% -----------------------------------------------------------------------------

Tabelle~\ref{tab:methodik-entscheidungen} fasst die zentralen
methodischen Entscheidungen zusammen.

\begin{table}[h]
    \centering
    \caption{Zentrale methodische Entscheidungen im Überblick}
    \label{tab:methodik-entscheidungen}
    \begin{tabular}{p{3.5cm} p{3.8cm} p{4.7cm}}
        \toprule
        \textbf{Entscheidung} & \textbf{Umsetzung} & \textbf{Begründung} \\
        \midrule
        Vergleichsdesign   & Paarvergleich mit/ohne xG           & Kausale Isolierung des xG-Effekts \\
        Datensatzvarianten & 12 Kombinationen                    & Robustheit über Vorverarbeitungsschritte \\
        Algorithmen        & 5 (LR, RF, LGBM, XGBoost, MLP)     & Breites Spektrum; faire Bewertung \\
        Tuning             & RandomizedSearchCV, StratifiedKFold & Vergleichbare Bedingungen für alle Modelle \\
        Datensplit         & 80:20, stratifiziert                & Repräsentative Klassenverteilung im Testset \\
        Leakage-Prävention & Shift(1) + Lag-1 für xG             & Temporal valide Vorhersagen \\
        Hauptmetrik        & F1 macro                            & Robust bei Klassen-Imbalance \cite{Powers2020} \\
        Sekundärmetrik     & F1 Draw (separat)                   & Fokus auf schwierigste Klasse \\
        \bottomrule
    \end{tabular}
\end{table}

Das folgende Kapitel~\ref{chap:modellierung} beschreibt die konkrete
Implementierung: Datengrundlage, Feature Engineering und Umsetzung
der ML-Pipeline.